\documentclass[10pt,a4paper]{amsart}
\usepackage{amsmath,amssymb,amsthm,graphicx}
% Retained preamble for the Stacks Project chapter snapshot in this directory.
% The chapter's own preamble.tex is a copy of the shared Stacks Project preamble; the block
% below is sliced MECHANICALLY, line for line, from it. Nothing was added, removed or
% redefined; the wrapper class, the xr-hyper/hyperref wiring and the externaldocument list
% are omitted because the assembler supplies the document class.
%
% lines taken from preamble.tex: 3, 5, 9, 11, 12, 156, 157, 158, ... (51 lines total)
%
% package coverage restored from the chapter's own preamble (2026-10-13): xy, xy-2cell

\usepackage[all]{xy}
\xyoption{2cell}
\UseAllTwocells

\newenvironment{reference}{\comment}{\endcomment}
\newenvironment{slogan}{\comment}{\endcomment}
\newenvironment{history}{\comment}{\endcomment}
\theoremstyle{plain}
\newtheorem{theorem}[subsection]{Theorem}
\newtheorem{proposition}[subsection]{Proposition}
\newtheorem{lemma}[subsection]{Lemma}
\theoremstyle{definition}
\newtheorem{definition}[subsection]{Definition}
\newtheorem{example}[subsection]{Example}
\newtheorem{exercise}[subsection]{Exercise}
\newtheorem{situation}[subsection]{Situation}
\theoremstyle{remark}
\newtheorem{remark}[subsection]{Remark}
\newtheorem{remarks}[subsection]{Remarks}
\numberwithin{equation}{subsection}
\def\lim{\mathop{\mathrm{lim}}\nolimits}
\def\colim{\mathop{\mathrm{colim}}\nolimits}
\def\Spec{\mathop{\mathrm{Spec}}}
\def\Hom{\mathop{\mathrm{Hom}}\nolimits}
\def\Ext{\mathop{\mathrm{Ext}}\nolimits}
\def\SheafHom{\mathop{\mathcal{H}\!\mathit{om}}\nolimits}
\def\SheafExt{\mathop{\mathcal{E}\!\mathit{xt}}\nolimits}
\def\Sch{\mathit{Sch}}
\def\Mor{\mathop{\mathrm{Mor}}\nolimits}
\def\Ob{\mathop{\mathrm{Ob}}\nolimits}
\def\Sh{\mathop{\mathit{Sh}}\nolimits}
\def\NL{\mathop{N\!L}\nolimits}
\def\CH{\mathop{\mathrm{CH}}\nolimits}
\def\proetale{{pro\text{-}\acute{e}tale}}
\def\etale{{\acute{e}tale}}
\def\QCoh{\mathit{QCoh}}
\def\Ker{\mathop{\mathrm{Ker}}}
\def\Im{\mathop{\mathrm{Im}}}
\def\Coker{\mathop{\mathrm{Coker}}}
\def\Coim{\mathop{\mathrm{Coim}}}
\DeclareMathSymbol{\boxtimes}{\mathbin}{AMSa}{"02}
\def\QCohstack{\mathcal{QC}\!\mathit{oh}}
\def\Cohstack{\mathcal{C}\!\mathit{oh}}
\def\Spacesstack{\mathcal{S}\!\mathit{paces}}
\def\Quotfunctor{\mathrm{Quot}}
\def\Hilbfunctor{\mathrm{Hilb}}
\def\Curvesstack{\mathcal{C}\!\mathit{urves}}
\def\Polarizedstack{\mathcal{P}\!\mathit{olarized}}
\def\Complexesstack{\mathcal{C}\!\mathit{omplexes}}
\def\Pic{\mathop{\mathrm{Pic}}\nolimits}
\def\Picardstack{\mathcal{P}\!\mathit{ic}}
\def\Picardfunctor{\mathrm{Pic}}
\def\Deformationcategory{\mathcal{D}\!\mathit{ef}}
% lane-s11: discard bibliographic \begin{reference} blocks entirely (incl. any \cite inside).
\newsavebox{\limrefbox}
\ifcsname reference\endcsname\expandafter\let\csname reference\endcsname\undefined\fi\ifcsname endreference\endcsname\expandafter\let\csname endreference\endcsname\undefined\fi
\def\citation#1{\relax}
\newenvironment{reference}{\setbox\limrefbox=\vbox\bgroup\let\cite\relax}{\egroup\ignorespaces}
\usepackage{cleveref}
\crefname{originallemma}{original lemma}{original lemmas}
\Crefname{originallemma}{Original lemma}{Original lemmas}
\makeatletter
\newlabel{lemma-directed-inverse-system-has-limit}{{01YX}{source}{}{sourcepointer-1}{}}
\newlabel{lemma-directed-inverse-system-has-limit@cref}{{[originallemma][0][]01YX}{[1][0][]source}{}{}{}}
\makeatother
\crefname{originalreference}{original reference}{original references}
\Crefname{originalreference}{Original reference}{Original references}
\makeatletter
\newlabel{topology-lemma-characterize-limit}{{0A2Q}{source}{}{sourcepointer-2}{}}
\newlabel{topology-lemma-characterize-limit@cref}{{[originalreference][0][]0A2Q}{[1][0][]source}{}{}{}}
\makeatother
\crefname{originallemma}{original lemma}{original lemmas}
\Crefname{originallemma}{Original lemma}{Original lemmas}
\makeatletter
\newlabel{lemma-inverse-limit-sets}{{0CUE}{source}{}{sourcepointer-3}{}}
\newlabel{lemma-inverse-limit-sets@cref}{{[originallemma][0][]0CUE}{[1][0][]source}{}{}{}}
\makeatother
\crefname{originallemma}{original lemma}{original lemmas}
\Crefname{originallemma}{Original lemma}{Original lemmas}
\makeatletter
\newlabel{lemma-directed-inverse-system-affine-schemes-has-limit}{{01YW}{source}{}{sourcepointer-4}{}}
\newlabel{lemma-directed-inverse-system-affine-schemes-has-limit@cref}{{[originallemma][0][]01YW}{[1][0][]source}{}{}{}}
\makeatother
\title{Topology of the inverse limit: three equivalent descriptions}
\author{The Stacks Project Authors}
\date{stacks-project:limits.tex:local-snapshot-20261012; GNU FDL 1.2 or later}
\begin{document}
\maketitle
\noindent\textit{Editorial note.} No omissions, hints or 'immediate' steps in the counted span.\par\medskip
\noindent\textit{Allowed standard prerequisites (editorial).}\begin{itemize}
\item Opens of the limit. Limits, Lemma lemma-descend-opens (same chapter)
\end{itemize}\medskip
\ifdefined\hypertarget\hypertarget{sourcepointer-1}{}\fi
\noindent\textit{Source reference.} Original lemma 01YX: sources/stacks-chow/tags.txt, source lines 2572--2572. Cross-chapter standard prerequisite cited by the author inside the counted span; it is used as a pointer and its own proof is not part of this item.\par\medskip
\ifdefined\hypertarget\hypertarget{sourcepointer-2}{}\fi
\noindent\textit{Source reference.} Original reference 0A2Q: sources/stacks-chow/tags.txt, source lines 12713--12713. Cross-chapter standard prerequisite cited by the author inside the counted span; it is used as a pointer and its own proof is not part of this item.\par\medskip
\ifdefined\hypertarget\hypertarget{sourcepointer-3}{}\fi
\noindent\textit{Source reference.} Original lemma 0CUE: sources/stacks-chow/tags.txt, source lines 16099--16099. Cross-chapter standard prerequisite cited by the author inside the counted span; it is used as a pointer and its own proof is not part of this item.\par\medskip
\ifdefined\hypertarget\hypertarget{sourcepointer-4}{}\fi
\noindent\textit{Source reference.} Original lemma 01YW: sources/stacks-chow/tags.txt, source lines 2571--2571. Cross-chapter standard prerequisite cited by the author inside the counted span; it is used as a pointer and its own proof is not part of this item.\par\medskip
\begin{lemma}
\label{lemma-inverse-limit-top}
\begin{reference}
\cite[IV, Proposition 8.2.9]{EGA}
\end{reference}
Let $S = \lim S_i$ be the limit of a directed inverse system
of schemes with affine transition morphisms
(Lemma \ref{lemma-directed-inverse-system-has-limit}). Then
$S_{top} = \lim_i S_{i, top}$ where $S_{top}$
indicates the underlying topological space of the scheme $S$.
\end{lemma}

\begin{proof}
We will use the criterion of
Topology, Lemma \ref{topology-lemma-characterize-limit}.
We have seen that $S_{set} = \lim_i S_{i, set}$ in
Lemma \ref{lemma-inverse-limit-sets}.
The maps $f_i : S \to S_i$ are morphisms of schemes
hence continuous. Thus $f_i^{-1}(U_i)$ is open for each
open $U_i \subset S_i$. Finally, 
let $s \in S$ and let $s \in V \subset S$ be an open
neighbourhood. Choose $0 \in I$ and choose an
affine open neighbourhood $U_0 \subset S_0$ of the image of $s$.
Then $f_0^{-1}(U_0) = \lim_{i \geq 0} f_{i0}^{-1}(U_0)$, see
Lemma \ref{lemma-directed-inverse-system-has-limit}.
Then $f_0^{-1}(U_0)$ and $f_{i0}^{-1}(U_0)$ are affine and
$$
\mathcal{O}_S(f_0^{-1}(U_0)) =
\colim_{i \geq 0} \mathcal{O}_{S_i}(f_{i0}^{-1}(U_0))
$$
either by the proof of
Lemma \ref{lemma-directed-inverse-system-has-limit}
or by
Lemma \ref{lemma-directed-inverse-system-affine-schemes-has-limit}.
Choose $a \in \mathcal{O}_S(f_0^{-1}(U_0))$ such that
$s \in D(a) \subset V$. This is possible because the principal
opens form a basis for the topology on the affine scheme $f_0^{-1}(U_0)$.
Then we can pick an $i \geq 0$ and
$a_i \in \mathcal{O}_{S_i}(f_{i0}^{-1}(U_0))$ mapping to $a$.
It follows that $D(a_i) \subset f_{i0}^{-1}(U_0) \subset S_i$
is an open subset whose inverse image in $S$ is $D(a)$.
This finishes the proof.
\end{proof}

\end{document}
